\section{Introduction}

Cancer cell lines are indispensable tools in translational cancer research, enabling mechanistic investigation, high-throughput drug screening, and functional validation of oncogenic processes. Their value as preclinical models, however, depends on how faithfully their transcriptional programmes reflect those of patient tumours. This correspondence cannot be assumed from tissue-of-origin labels alone: cell lines can diverge from primary tumours during prolonged \textit{in vitro} culture and through the absence of the tumour microenvironment \citep{Bose2024,Mirabelli2019CancerCellLines,BenDavid2018}. Patient tumours also show substantial transcriptional heterogeneity within and across cancer types \citep{PanCancerWholeGenome2020}. Model selection is therefore a bidirectional transcriptomic similarity-ranking problem: patient tumours can be used to rank candidate cell-line models, and cell lines can be used to rank patient tumours according to transcriptomic proximity.

Several computational approaches have addressed related aspects of matching tumours to cell lines. Celligner performs global transcriptomic alignment between tumours and cell lines \citep{Warren2021GlobalAlignment}, MFmap uses semi-supervised generative mapping to match cell lines to tumours and cancer subtypes \citep{Zhang2021MFmap}, and OncoNPC predicts tumour origin from molecular features \citep{Moon2023CUP}. Together, these approaches have advanced molecular mapping, tumour classification, and model-selection strategies. However, they do not directly provide a bidirectional framework for resolving stable neighbourhoods within a defined cell-line panel, identifying isolated and bridge-like models, or evaluating feature--distance sensitivity. They also do not allow candidate cell-line models to be ranked against individual patient tumours, nor patient tumours to be ranked from each cell-line profile. This leaves a practical gap: researchers need to know not only whether a cancer type is represented in a cell-line panel, but which specific models best match the patient tumour contexts relevant to their experiment.

Transcriptomic matching between patient tumours and cell lines also faces two structural sources of bias. First, bulk RNA-seq profiles from patient tumours combine malignant, stromal, endothelial, and immune-derived signals, whereas cancer cell lines are malignant-cell-enriched systems. Without tumour-purity filtering, similarity scores may partly reflect microenvironmental admixture rather than tumour-intrinsic transcriptional programmes \citep{Yoshihara2013}. Second, transcriptomic proximity is sensitive to the feature space and distance geometry used to compare samples. A single gene-selection criterion or distance metric may emphasise only one aspect of tumour biology, making model rankings vulnerable to analytical choices or lineage-specific expression structure \citep{Solorio-Fernandez2020,Kallberg2021}. Technical variation between sequencing datasets can introduce batch effects and other unwanted variation that further distort transcriptomic proximity \citep{Leek2010BatchEffects,Leek2012}.

Together, these requirements motivate a workflow that combines biological filtering with analytical redundancy. Tumour-purity filtering and source harmonisation make patient tumours and cell lines more comparable before similarity estimation. Multi-view feature selection evaluates feature--distance representations spanning dispersion, entropy, principal component structure, co-expression connectivity, and complementary distance metrics \citep{Solorio-Fernandez2020,Kallberg2021}. Within this representation set, weighted gene co-expression network analysis contributes co-expression connectivity features \citep{Langfelder2008WGCNA}. Consensus clustering identifies cluster structure that persists under resampling \citep{Monti2003ConsensusClustering}. \texttt{ConsensusClusterPlus} provided the consensus-clustering implementation used in this workflow \citep{Wilkerson2010ConsensusClusterPlus}, and adaptive tumour-neighbourhood analysis identified associations that persist across representations. Graph-based resolution uses global-best and local-best neighbour sets to define interpretable model structures, including stable components, isolated models, and bridge-like cell lines. Differential expression analysis identifies marker genes for the resolved groups \citep{Love2014}. Functional enrichment analysis provides biological context for those marker sets \citep{Raudvere2019gProfiler}. Similarity-based assessment by nearest-neighbour lineage matching quantified whether the resulting feature space preserved the expected tumour-lineage relationships.

To address this gap, this study applied a bidirectional transcriptomic framework to cancer cell lines from the Deutsche Sammlung von Mikroorganismen und Zellkulturen (DSMZ) panel and clinical tumour RNA-seq cohorts. The analysis covers breast cancer, neuroblastoma, and retinoblastoma cohorts, which provide distinct test cases spanning a heterogeneous adult epithelial cancer and two paediatric developmental tumour lineages. Specifically, the study aims to (i) evaluate bidirectional transcriptomic similarity ranking between patient tumours and cell-line expression profiles, (ii) define stable model neighbourhoods that persist across feature--distance choices, and (iii) identify transcriptomically atypical models whose patient-tumour similarity profiles are not fully explained by lineage annotations.
